
\begin{table}[htbp]
    \centering
    \caption{Exploratory Factor Analysis Model Selection Statistics.}
    \label{tab:EFA-10}
    \begin{threeparttable}
        \begin{tabular}{cccccc
        }
            \toprule
            {# Factors} & {PCA Eigenvalue} & {AIC} & {BIC} & {CFI} & {RMSEA}\\
            \midrule
            1 & 3.239 & 3734.880 & 3793.425 & 0.463 & 0.220\\
            2 & 1.545 & 3630.132 & 3715.022 & 0.724 & 0.183\\
            3 & 1.261 & 3554.424 & {\textbf{3662.733}} & 0.917 & 0.121\\
            4 & \textbf{1.183} & 3537.297 & 3666.097 & 0.972 & 0.089\\
            5 & 0.771 & {\textbf{3527.484}} & 3673.847 & 1.000 & 0.000\\
            6 & 0.682 & 3536.200 & 3697.198 & 1.000 & 0.000\\
            \bottomrule
        \end{tabular}
        \vspace{5pt}
        \begin{tablenotes}
            \small
            \item Note: Boldface typeset indicates number of selected factors.  These are 4 by eigenvalues, 5 by AIC and 3 by BIC.
        \end{tablenotes}
    \end{threeparttable}
\end{table}


\begin{table}[htbp]
    \centering
    \caption{Exploratory Factor Analysis Loadings}
    \label{tab:Fac-10}
    \begin{threeparttable}
        \begin{tabular}{l|ccccc|cc}
            \toprule
            Task-Version
            & \multicolumn{5}{c|}{Factor Loadings} & \multicolumn{2}{c}{Proportion of Variance} \\
            \cmidrule{2-8}
             & 1 & 2 & 3 & 4 & 5 & Unique & Shared \\
            \midrule
            Brentano v1 & \textit{0.401} & 0.145 & \textit{0.279} & 0.124 & -0.119 & 0.711 & 0.289 \\
            Brentano v2 & \textit{0.315} & 0.008 & \textit{0.336} & 0.019 & 0.041 & 0.785 & 0.215 \\
            Ebbinghaus v1 & 0.096 & 0.014 & 0.124 & 0.223 & \textbf{0.724} & 0.401 & 0.599 \\
            Ebbinghaus v2 & 0.163 & 0.043 & 0.103 & \textbf{0.854} & 0.267 & 0.160 & 0.840 \\
            Poggendorf v1 & 0.119 & 0.072 & \textbf{0.904} & 0.104 & 0.013 & 0.152 & 0.848 \\
            Poggendorf v2 & 0.154 & 0.101 & \textbf{0.704} & 0.021 & 0.155 & 0.446 & 0.554 \\
            Ponzo v1 & 0.189 & \textbf{0.811} & 0.092 & 0.175 & -0.115 & 0.254 & 0.746 \\
            Ponzo v2 & 0.072 & \textbf{0.976} & 0.094 & -0.113 & 0.144 & 0.000 & 1.000 \\
            Zoellner v1 & \textbf{0.870} & 0.147 & 0.034 & 0.080 & 0.181 & 0.182 & 0.818 \\
            Zoellner v2 & \textbf{0.771} & 0.070 & 0.209 & 0.081 & 0.049 & 0.348 & 0.652 \\
            \bottomrule
        \end{tabular}
        \vspace{5pt}
        \begin{tablenotes}
            \small
            \item Note. Boldfaced loadings highlight the relations between factors and tasks.
        \end{tablenotes}
    \end{threeparttable}
\end{table}


\begin{table}[ht]
\centering
\caption{Factor Loadings For CFA Model of Visual Illusions.}
\label{tab:loadings-new}
\begin{tabular}{lccccccccl}
\toprule
\textbf{Task} & \multicolumn{5}{c}{\textbf{Nuisance Factors}} & \multicolumn{2}{c}{\textbf{Exploratory Factors}} & \textbf{Unique Variance} \\
\cmidrule(lr){2-6} \cmidrule(lr){7-8}
 & F1 & F2 & F3 & F4 & F5 & F6 & F7 & \\
\midrule
Brentano 1   & 0.270 & --    & --    & --    & --    & 0.569 & -0.073 & 0.477 \\
Brentano 2   & -0.075 & --   & --    & --    & --    & 0.494 & -0.007 & 0.539 \\
Ebbinghaus 1 & --    & 0.621 & --    & --    & --    & 0.251 &  0.160 & 0.246 \\
Ebbinghaus 2 & --    & 0.625 & --    & --    & --    & 0.354 & -0.040 & 0.249 \\
Poggendorf 1 & --    & --    & 0.623 & --    & --    & 0.528 & -0.012 & 0.393 \\
Poggendorf 2 & --    & --    & 0.565 & --    & --    & 0.508 &  0.066 & 0.342 \\
Ponzo 1      & --    & --    & --    & 0.824 & --    & 0.352 & -0.011 & 0.194 \\
Ponzo 2      & --    & --    & --    & 0.842 & --    & 0.272 &  0.110 & 0.165 \\
Zöllner 1    & --    & --    & --    & --    & 0.588 & 0.574 &  0.011 & 0.321 \\
Zöllner 2    & --    & --    & --    & --    & 0.555 & 0.616 & -0.042 & 0.298 \\
\midrule
\textbf{Prop. Var.} & 0.021 & 0.081 & 0.078 & 0.141 & 0.069 & 0.233 & 0.055 & 0.322 \\
\bottomrule
\end{tabular}

\vspace{0.5em}
\raggedright
\footnotesize
\textit{Note.} We fit a 7-factor hierarchical Bayesian model. The first five factors were used to capture within-version idiosyncrasies, with all other loadings zeroed out for other tasks. The last two factors were general and shared across all versions and tasks. All loadings were scaled to reflect variance decomposition in the correlation metric, such that each value represents a factor’s contribution to shared variance. The first general factor explained over four times as much variance in the correlation structure as the second. For more details, refer to \textcite{Rouder.etal.2024a}.
\end{table}
